% Generated draft: requires mathematical and source review.
\documentclass[11pt]{article}
\usepackage[margin=25mm]{geometry}
\usepackage{fontspec}
\setmainfont{Latin Modern Roman}
\usepackage{amsmath,amssymb,mathtools,unicode-math}
\setmathfont{Latin Modern Math}
\usepackage{enumitem,needspace,xurl,hyperref}
\setlength{\parindent}{0pt}
\setlength{\parskip}{5pt}
\setlength{\emergencystretch}{3em}
\newcommand{\N}{\mathbb{N}}
\newcommand{\Z}{\mathbb{Z}}
\newcommand{\Q}{\mathbb{Q}}
\newcommand{\R}{\mathbb{R}}
\newcommand{\C}{\mathbb{C}}
\newcommand{\F}{\mathbb{F}}
\newcommand{\A}{\mathbb{A}}
\newcommand{\PP}{\mathbb{P}}
\newcommand{\E}{\mathbb{E}}
\newcommand{\eps}{\varepsilon}
\newcommand{\dd}{\,\mathrm{d}}
\begin{document}
\section*{UnsolvedMath ID: 30002767}
\textbf{Source problem code:} catalogue entry
\subsection{Definitions and mathematical statement}
For a group \(G\), define its integral cohomological dimension by
\[
\operatorname{cd}_{\Z}(G)=\operatorname{pd}_{\Z G}\Z,
\]
where \(\Z\) is the trivial \(\Z G\)-module. Its geometric dimension \(\operatorname{gd}(G)\) is the least dimension of a CW complex \(K(G,1)\), meaning a connected CW complex with fundamental group \(G\) and contractible universal cover. The Eilenberg--Ganea conjecture in the case at issue asserts that \(\operatorname{cd}_{\Z}(G)=2\) implies \(\operatorname{gd}(G)=2\).

A simplicial complex is flag if every finite set of pairwise adjacent vertices spans a simplex. For a finite flag complex \(L\), let \(V(L)\) and \(E(L)\) denote its vertex and edge sets, and define the associated right-angled Coxeter group by
\[
W_L=\left\langle s_v\ (v\in V(L))\ \middle|\ s_v^2=1\ (v\in V(L)),\quad s_vs_w=s_ws_v\ (\{v,w\}\in E(L))\right\rangle.
\]
A subgroup \(H\leq W_L\) has finite index if \([W_L:H]<\infty\), and is torsion-free if its identity is its only element of finite order.

For a specified coefficient ring \(A\), a nonempty complex \(L\) is acyclic over \(A\) if \(\widetilde H_k(L;A)=0\) for every integer \(k\geq 0\). The supplied record does not specify the coefficients intended by ``acyclic.'' A unitary representation of \(\pi_1(L)\) is a homomorphism into the unitary group of a complex Hilbert space; it is nontrivial if some group element acts other than as the identity. The record does not specify whether the representation must be finite-dimensional.

\textbf{Supplied question, subject to these unresolved conventions.} Let \(L\) be a finite acyclic flag complex of dimension exactly two whose fundamental group admits a nontrivial unitary representation. Let \(H\leq W_L\) be any torsion-free subgroup of finite index. Taking
\[
\operatorname{cd}_{\Z}(H)=2
\]
as the premise asserted in the supplied record, must every such \(L\) and \(H\) satisfy
\[
\operatorname{gd}(H)=3?
\]
An affirmative answer would make these groups counterexamples to the stated Eilenberg--Ganea conjecture. The original wording restricts attention to groups arising from the Bestvina--Davis construction; the supplied clean formulation describes them using the hypotheses above. Whether that description includes all necessary construction hypotheses requires source review. No coefficient convention, representation dimension, or additional construction hypothesis is resolved by this draft.
\subsection{Short English statement}
The supplied question asks whether every torsion-free finite-index subgroup in the stated Bestvina--Davis family has geometric dimension three despite having integral cohomological dimension two. The meaning of acyclicity, the allowed dimension of the unitary representation, and the completeness of the construction hypotheses require review.
\subsection{Sources}
\begin{itemize}
\item Supplied local UnsolvedMath registry record, permanent ID \(30002767\), titled ``Bestvina--Davis Candidates for the Eilenberg--Ganea Conjecture,'' source problem code \texttt{catalogue entry}. The record belongs to the Oberwolfach Reports Open Problems collection and gives proposed year 2015; it names no individual proposer.
\item Supplied report citation: \emph{Geometric Topology} (2015), Oberwolfach Reports, DOI \texttt{10.4171/owr/2015/3}, pp.\ 187--233. \url{https://doi.org/10.4171/owr/2015/3}. The background describes these reports as records of mathematical workshops and their open-problem discussions.
\item Supplied literature source: Bestvina--Brady, \emph{Inventiones Mathematicae}. \url{https://people.math.osu.edu/davis.12/courses/8800-20/Bestvina-Brady.pdf}.
\item Supplied literature source: Davis, Eilenberg--Ganea survey. \url{https://people.math.osu.edu/davis.12/np1.pdf}.
\item Literature assessment dated exactly \texttt{2026-08-21}: the supplied curation labels the problem open and states that the Eilenberg--Ganea question for torsion-free groups of cohomological dimension two remains open, with the Bestvina--Davis construction providing candidates rather than a definitive resolution. This is supplied dated curation, not an independently verified literature assessment; the background explicitly requests independent verification of time-sensitive status and cited claims.
\item The supplied record labels its statement \texttt{corrected\_verified} and says that it was checked against the official Oberwolfach report and preserves the report's hypotheses, notation, and target. These are supplied verification claims, not verification performed for this draft.
\end{itemize}
\end{document}
